% =============================================================================
% CHƯƠNG: AI PREDICTION VỚI EXPLAINABLE AI (XAI)
% =============================================================================

\section{AI PREDICTION VỚI EXPLAINABLE AI (XAI)}

\subsection{Bối cảnh}

Trong hệ thống phân tích tiền mã hóa, ngoài việc thu thập tin tức và phân tích sentiment, yêu cầu quan trọng là có thể:
\begin{itemize}
    \item Dự đoán xu hướng giá (UP/DOWN) trong các khung thời gian khác nhau (1h, 4h, 24h)
    \item Kết hợp dữ liệu tin tức với giá lịch sử để đưa ra dự đoán chính xác hơn
    \item \textbf{Giải thích lý do} tại sao mô hình đưa ra dự đoán đó (Explainable AI)
\end{itemize}

\subsection{Vấn đề với AI "Hộp đen"}

\subsubsection{Tại sao cần Explainable AI?}

\begin{enumerate}
    \item \textbf{Thiếu tin tưởng}: Người dùng không muốn đầu tư dựa trên dự đoán mà không hiểu lý do
    \item \textbf{Regulatory compliance}: Các quy định tài chính yêu cầu giải thích quyết định AI
    \item \textbf{Debug và cải thiện}: Không thể cải thiện model nếu không hiểu nó hoạt động như thế nào
    \item \textbf{Bias detection}: Phát hiện các yếu tố không mong muốn ảnh hưởng đến dự đoán
\end{enumerate}

\subsection{Giải pháp: XGBoost + SHAP Pipeline}

\subsubsection{Kiến trúc Prediction Pipeline}

\begin{figure}[H]
\centering
\begin{tikzpicture}[
    node distance=1.2cm,
    box/.style={rectangle, draw, minimum width=2.5cm, minimum height=0.8cm, align=center, font=\small, rounded corners}
]

% Data Sources
\node[box, fill=yellow!30] (prices) {Price Data\\Binance 1h};
\node[box, fill=cyan!30, right=2cm of prices] (news) {News Data\\Sentiment Scores};

% Feature Engineering
\node[box, fill=orange!30, below=1.5cm of prices] (techfeat) {Technical Features\\RSI, MACD, SMA};
\node[box, fill=orange!30, below=1.5cm of news] (newsfeat) {News Features\\Sentiment MA};

% Model
\node[box, fill=green!30, below right=1.5cm and -0.5cm of techfeat] (xgboost) {XGBoost\\Regressor};

% Explainer
\node[box, fill=red!30, below=1.2cm of xgboost] (shap) {SHAP\\Explainer};

% Outputs
\node[box, fill=purple!30, below left=1.2cm and 0.5cm of shap] (prediction) {Predicted Price\\+ Confidence};
\node[box, fill=purple!30, below right=1.2cm and 0.5cm of shap] (explanation) {Feature\\Contributions};

% Arrows
\draw[->, thick] (prices) -- (techfeat);
\draw[->, thick] (news) -- (newsfeat);
\draw[->, thick] (techfeat) -- (xgboost);
\draw[->, thick] (newsfeat) -- (xgboost);
\draw[->, thick] (xgboost) -- (shap);
\draw[->, thick] (shap) -- (prediction);
\draw[->, thick] (shap) -- (explanation);

\end{tikzpicture}
\caption{AI Prediction Pipeline với SHAP Explainer}
\end{figure}

\subsubsection{Prediction Pipeline Implementation}

\begin{lstlisting}[language=Python, caption=Prediction Pipeline (ai-service-python/internal/prediction.py)]
class PredictionPipeline:
    """Prediction pipeline with XGBoost and SHAP"""
    
    def __init__(self, analyzer: Optional[Analyzer] = None):
        self.models = {}  # Store models for each symbol+horizon
        self.feature_configs = {}  # Store feature names
        self.explainers = {}  # Store SHAP explainers
        self.analyzer = analyzer  # Sentiment analyzer

    def load_price_data(self, symbol: str = "BTCUSDT", years: int = 2):
        """Load price data from database"""
        query = """
            SELECT symbol, time, interval, open, high, low, close, volume
            FROM market_prices
            WHERE symbol = %s AND interval = '1h'
            AND time >= NOW() - INTERVAL '%s years'
            ORDER BY time ASC
        """
        return pd.read_sql(query, db.conn, params=(symbol, years))
    
    def load_news_data(self, days: int = 730):
        """Load news with sentiment scores"""
        query = """
            SELECT id, source_id, title, content_text, published_at,
                   sentiment_score, url, language
            FROM articles
            WHERE crawled_at >= NOW() - INTERVAL '%s days'
            ORDER BY crawled_at ASC
        """
        return pd.read_sql(query, db.conn, params=(days,))
\end{lstlisting}

\subsection{Feature Engineering}

\subsubsection{Technical Indicators}

\begin{lstlisting}[language=Python, caption=Price Feature Engineering]
def create_price_features(self, df: pd.DataFrame) -> pd.DataFrame:
    """Create technical indicator features"""
    df = df.copy().sort_values('time').reset_index(drop=True)
    
    # Returns
    df['return_1h'] = df['close'].pct_change(1)
    df['return_24h'] = df['close'].pct_change(24)
    
    # Volatility
    df['volatility_24h'] = df['return_1h'].rolling(24).std()
    
    # Simple Moving Average
    df['sma_7'] = df['close'].rolling(7).mean()
    df['sma_24'] = df['close'].rolling(24).mean()
    
    # RSI (Relative Strength Index)
    delta = df['close'].diff()
    gain = delta.where(delta > 0, 0).rolling(14).mean()
    loss = -delta.where(delta < 0, 0).rolling(14).mean()
    rs = gain / loss
    df['rsi'] = 100 - (100 / (1 + rs))
    
    # MACD (Moving Average Convergence Divergence)
    ema12 = df['close'].ewm(span=12).mean()
    ema26 = df['close'].ewm(span=26).mean()
    df['macd'] = ema12 - ema26
    
    # Lag features
    df['close_lag_1'] = df['close'].shift(1)
    df['close_lag_24'] = df['close'].shift(24)
    
    return df
\end{lstlisting}

\subsubsection{News Sentiment Features}

\begin{lstlisting}[language=Python, caption=News Feature Engineering with Time Alignment]
def process_news(self, news_df: pd.DataFrame, price_df: pd.DataFrame):
    """Process news and align with price time buckets"""
    
    # Round to hourly buckets for alignment
    news_df['time_bucket'] = news_df['time'].dt.floor('1h')
    price_df['time_bucket'] = pd.to_datetime(price_df['time']).dt.floor('1h')
    
    # Aggregate news sentiment per hour
    news_agg = news_df.groupby('time_bucket').agg({
        'sentiment': ['mean', 'count', 'std', lambda x: (x > 0.1).sum()]
    }).reset_index()
    news_agg.columns = ['time_bucket', 'sentiment_mean', 'news_count', 
                        'sentiment_std', 'positive_count']
    
    # Merge with price
    result = price_df.merge(news_agg, on='time_bucket', how='left')
    result['sentiment_mean'] = result['sentiment_mean'].fillna(0)
    result['news_count'] = result['news_count'].fillna(0)
    
    # Sentiment-derived features
    result['sentiment_positive_ratio'] = result['positive_count'] / (result['news_count'] + 1)
    
    # Lag features for sentiment
    result['sentiment_lag_1h'] = result['sentiment_mean'].shift(1).fillna(0)
    result['sentiment_lag_6h'] = result['sentiment_mean'].shift(6).fillna(0)
    
    # Rolling features
    result['sentiment_ma_24h'] = result['sentiment_mean'].rolling(24, min_periods=1).mean()
    result['sentiment_std_24h'] = result['sentiment_mean'].rolling(24, min_periods=1).std()
    
    return result
\end{lstlisting}

\subsection{Model Training với XGBoost}

\begin{lstlisting}[language=Python, caption=XGBoost Training with Feature Selection]
def train(self, symbol: str = "BTCUSDT", horizon_hours: int = 4, years: int = 2):
    """Train XGBoost model"""
    
    # Load and prepare data
    price_df = self.load_price_data(symbol, years)
    news_df = self.load_news_data(days=years*365)
    
    price_df = self.create_price_features(price_df)
    df = self.process_news(news_df, price_df)
    
    # Create target: future price after horizon_hours
    df['target'] = df['close'].shift(-horizon_hours)
    df = df.dropna(subset=['target'])
    
    # Feature selection: ensure balanced news vs technical
    feature_cols = self.select_features(df, df['target'], n=30)
    
    # Prepare data
    X = df[feature_cols].fillna(0)
    y = df['target']
    
    # Train/test split (80/20)
    split_idx = int(len(X) * 0.8)
    X_train, X_test = X.iloc[:split_idx], X.iloc[split_idx:]
    y_train, y_test = y.iloc[:split_idx], y.iloc[split_idx:]
    
    # Train XGBoost
    model = xgb.XGBRegressor(
        n_estimators=200, 
        max_depth=6, 
        learning_rate=0.05, 
        subsample=0.8,
        colsample_bytree=0.8,
        random_state=42
    )
    model.fit(X_train, y_train)
    
    # Initialize SHAP explainer
    self.explainers[model_key] = shap.TreeExplainer(model)
    
    # Evaluate
    y_pred = model.predict(X_test)
    rmse = np.sqrt(mean_squared_error(y_test, y_pred))
    
    # Directional accuracy (UP/DOWN prediction)
    direction_true = np.sign(y_test.diff().iloc[1:])
    direction_pred = np.sign(pd.Series(y_pred).diff().iloc[1:])
    dir_acc = (direction_true == direction_pred).mean() * 100
    
    return {'rmse': rmse, 'directional_accuracy': dir_acc}
\end{lstlisting}

\subsection{SHAP Explainability}

\subsubsection{Giới thiệu SHAP (SHapley Additive exPlanations)}

SHAP là phương pháp XAI dựa trên lý thuyết trò chơi Shapley value:
\begin{itemize}
    \item \textbf{Shapley Value}: Phân bổ "đóng góp" của mỗi feature vào prediction cuối cùng
    \item \textbf{Local Explanation}: Giải thích cho từng prediction cụ thể
    \item \textbf{Feature Importance}: Tổng hợp importance từ nhiều predictions
\end{itemize}

\textbf{Công thức Shapley Value}:
$$\phi_i = \sum_{S \subseteq N \setminus \{i\}} \frac{|S|!(|N|-|S|-1)!}{|N|!} [v(S \cup \{i\}) - v(S)]$$

Trong đó:
\begin{itemize}
    \item $\phi_i$: Shapley value của feature $i$
    \item $N$: Tập tất cả features
    \item $S$: Một subset của features (không bao gồm $i$)
    \item $v(S)$: Giá trị prediction khi chỉ sử dụng features trong $S$
\end{itemize}

\subsubsection{SHAP Implementation}

\begin{lstlisting}[language=Python, caption=SHAP Explanation trong Prediction]
def predict(self, symbol: str = "BTCUSDT", horizon_hours: int = 4) -> Dict:
    """Make prediction with SHAP explanation"""
    
    model = self.models.get(model_key)
    explainer = self.explainers.get(model_key)
    feature_names = self.feature_configs.get(model_key)
    
    # Get latest data
    X = latest[feature_names].fillna(0)
    
    # Make prediction
    current_price = float(latest['close'].iloc[0])
    predicted_price = float(model.predict(X)[0])
    
    # SHAP Explanation
    explanation = {}
    if explainer:
        shap_values = explainer.shap_values(X)
        feature_contrib = dict(zip(feature_names, shap_values[0]))
        
        # Sort by absolute contribution
        top_contrib = sorted(feature_contrib.items(), 
                           key=lambda x: abs(x[1]), reverse=True)[:10]
        explanation = {k: float(v) for k, v in top_contrib}
    
    # Separate News vs Technical features
    news_features = [f for f in feature_names 
                     if 'sentiment' in f.lower() or 'news' in f.lower()]
    tech_features = [f for f in feature_names if f not in news_features]
    
    return {
        "predicted_price": str(predicted_price),
        "current_price": str(current_price),
        "predicted_change_pct": str((predicted_price - current_price) / current_price * 100),
        
        # XAI: Feature analysis
        "feature_analysis": {
            "top_news_features": [...],
            "top_technical_indicators": [...],
            "shap_contributions": explanation  # SHAP values for each feature
        },
        
        # Human-readable explanation
        "explanation": {
            "primary_factors": {
                "most_important_feature": top_features[0][0],
                "is_news_important": bool(news dominates),
                "news_influence_pct": f"{news_pct:.1f}%"
            }
        }
    }
\end{lstlisting}

\subsection{Phân tích Nhân quả (Causal Analysis)}

\subsubsection{Kết hợp Tin tức và Giá}

Hệ thống thực hiện phân tích nhân quả bằng cách:

\begin{enumerate}
    \item \textbf{Time-aligned Features}: Tin tức được align với khung giờ giá
    \item \textbf{Lag Features}: Sentiment lag 1h, 6h để capture causal delay
    \item \textbf{SHAP Contributions}: Xác định feature nào thực sự ảnh hưởng đến prediction
\end{enumerate}

\begin{lstlisting}[language=Python, caption=Causal Feature Extraction]
# Lag features capture causal relationship
result['sentiment_lag_1h'] = result['sentiment_mean'].shift(1)
result['sentiment_lag_6h'] = result['sentiment_mean'].shift(6)

# Rolling features for longer-term trends
result['sentiment_ma_24h'] = result['sentiment_mean'].rolling(24).mean()
\end{lstlisting}

\subsubsection{Top News Articles với Keywords}

\begin{lstlisting}[language=Python, caption=Extract Top Influential News]
# Get recent news relevant to prediction
lookback_hours = max(horizon_hours * 2, 48)
cutoff_time = news_df['time'].max() - timedelta(hours=lookback_hours)
recent_news = news_df[news_df['time'] >= cutoff_time]

# Extract keywords using Attention weights
for idx, row in recent_news.head(3).iterrows():
    text = row['title'] + " " + row['content_text']
    keywords, _ = self.analyzer.analyze(text, lang=row['language'], top_k=5)
    
    top_news_articles.append({
        "id": row['id'],
        "title": row['title'],
        "sentiment_score": row['sentiment_score'],
        "keywords": keywords[:5]  # Top 5 keywords from attention
    })
\end{lstlisting}

\subsection{Frontend: Prediction Panel}

\begin{lstlisting}[language=TypeScript, caption=PredictionPanel Component (React)]
interface PredictionResult {
  predicted_price: string;
  current_price: string;
  predicted_change_pct: string;
  confidence_score: string;
  explanation: {
    primary_factors: {
      most_important_feature: string;
      is_news_important: boolean;
      news_influence_pct: string;
    }
  };
  top_news_articles: Array<{
    id: number;
    title: string;
    sentiment_score: number;
    keywords: string[];
  }>;
  feature_analysis: {
    top_technical_indicators: Array<{ feature: string; importance: number }>;
  };
}

// Display prediction with explanation
<div className="prediction-panel">
  <h3>AI Price Prediction - {symbol} • {horizonHours}h horizon</h3>
  
  {/* Price prediction */}
  <p>Current: ${currentPrice}</p>
  <p className={isPositive ? 'green' : 'red'}>
    Predicted: ${predictedPrice} ({changePct}%)
  </p>
  
  {/* Key factors - XAI */}
  <h4>Key Factors</h4>
  <p>Primary Factor: {prediction.explanation.primary_factors.most_important_feature}</p>
  <p>News Influence: {prediction.explanation.primary_factors.news_influence_pct}</p>
  
  {/* Top news articles */}
  <h4>Recent News ({prediction.metadata.recent_news_analyzed} analyzed)</h4>
  {prediction.top_news_articles.map(article => (
    <div key={article.id}>
      <p>{article.title}</p>
      <span>Sentiment: {article.sentiment_score}</span>
      <span>Keywords: {article.keywords.join(', ')}</span>
    </div>
  ))}
  
  {/* Technical indicators */}
  <h4>Top Technical Indicators</h4>
  {prediction.feature_analysis.top_technical_indicators.map(ind => (
    <span className="badge">{ind.feature}</span>
  ))}
</div>
\end{lstlisting}

\subsection{Confidence Score Calculation}

\begin{lstlisting}[language=Python, caption=Multi-factor Confidence Score]
# 1. Feature coverage (enough features?)
feature_coverage = len(feature_names) / 30.0

# 2. News availability (recent news available?)
news_availability = min(1.0, recent_news_count / max(10, horizon_hours // 4))

# 3. Balance score (model diversity?)
news_feat_ratio = len(news_features) / len(feature_names)
balance_score = 1.0 - abs(news_feat_ratio - 0.3)  # Ideal: 30% news features

# Combined confidence
confidence = (feature_coverage * 0.4 + news_availability * 0.3 + balance_score * 0.3)
\end{lstlisting}

\textbf{Công thức Confidence}:
$$confidence = 0.4 \times coverage + 0.3 \times news\_availability + 0.3 \times balance$$

\subsection{Kết luận}

\begin{table}[H]
\centering
\caption{So sánh Black-box vs Explainable AI}
\begin{tabular}{|p{4cm}|p{5cm}|p{5cm}|}
\hline
\textbf{Tiêu chí} & \textbf{Black-box AI} & \textbf{XAI (SHAP)} \\
\hline
Độ tin cậy người dùng & Thấp & Cao \\
\hline
Debug capability & Khó & Dễ - thấy feature contributions \\
\hline
Regulatory compliance & Không đạt & Đạt chuẩn \\
\hline
Bias detection & Không thể & Có thể phát hiện \\
\hline
Complexity & Thấp & Cao hơn (cần SHAP computation) \\
\hline
\end{tabular}
\end{table}

\textbf{Điểm mạnh của giải pháp}:
\begin{itemize}
    \item Kết hợp cả Technical Indicators và News Sentiment
    \item SHAP explanation cho từng prediction
    \item Hiển thị top news articles với keywords để người dùng hiểu lý do
    \item Confidence score dựa trên nhiều yếu tố
    \item Support multiple prediction horizons (1h, 4h, 24h)
\end{itemize}

\newpage
